\chapter{操作可行性}

\section{多端响应式与国际化支持}
LingoBridge 自立项以来便强调极简的用户交互路径，以适应华文教育中不同年龄段、不同国别师生的教学使用门槛：
\begin{itemize}
  \item \textbf{高美学响应式 UI}：前端界面遵循精美的美学排版标准，从大屏 of 课程管理 Tab 到平板的在线白板、手机端的日程卡片，均能实现 100\% 像素级自适应适配。
  \item \textbf{四语无感无偏翻译}：为了支持跨国交流，系统全局集成了中、英、俄、哈四语的实时国际化切换（多语资源存储于本地 `LanguageContext`），避免了跨国师生因语言屏障引发的认知障碍。
\end{itemize}

\section{双角色交互门槛分析}
\begin{itemize}
  \item \textbf{教师端（零复杂度运维）}：教师可使用图形化排课日历进行秒级新增/编辑，并可通过 \textbf{Standard Excel 作业表格} 一键上传。Excel 的解析逻辑强绑定课时节点（`lessonNodeId`），其去重校验（Upsert Whitelist 机制）基于唯一索引，即使反复上传也绝不覆盖已有批改结果，大大减轻了备课心智负担。
  \item \textbf{学生端（即入即学）}：学生登录后可直观进入 active 课堂，页面包含微脉冲“双雷达”网络自连接雷达动画。在线跟读作业时，界面提供 \textbf{3-slot 自动磁贴式录音音轨对比}，发音评分系统自动高亮发音盲区（基于 ASR 云打分），实现了从“听-读-纠”的超简操作闭环。
\end{itemize}
